\documentclass[11pt,a4paper]{article}

\usepackage[utf8]{inputenc}
\usepackage[T1]{fontenc}
\usepackage{lmodern}
\usepackage{microtype}
\usepackage[english]{babel}
\usepackage[a4paper,margin=2.5cm,headheight=14pt]{geometry}
\usepackage{xcolor}
\usepackage{enumitem}
\usepackage{booktabs}
\usepackage{fancyhdr}
\usepackage{hyperref}
\usepackage{titlesec}

\hypersetup{
    colorlinks=true,
    linkcolor=black,
    urlcolor=blue!60!black,
}

\titleformat{\section}{\large\bfseries}{\thesection.}{0.6em}{}

\pagestyle{fancy}
\fancyhf{}
\fancyhead[L]{CSE 341 -- Recipix v4.1}
\fancyhead[R]{D5 (Part 2)}
\fancyfoot[C]{\thepage}

\title{\textbf{D5 -- Retrospective \& Self-Assessment (Part 2)} \\
       \large Recipix v4.1 -- CSE 341 Semester Project}
\author{Berk Hakan \"Oge (210104004132)}
\date{23 May 2026}

\begin{document}
\maketitle
\thispagestyle{fancy}


\section{Retrospective (\textonehalf{} page)}

\paragraph{What I would change if I started over.}
I would lock the Part-2 runtime contract on day one of the project,
not on day one of Part~2. The \texttt{tokens.py} agreement was what
made Part~1 work as parallel code with İsmail, and the equivalent for
Part~2 -- the runtime contract document with the in-place AST
annotation rule, the runtime value dataclasses, the action-verb
signature table, the environment chain, and the agreed-upon expected
outputs for samples~1 and~2 -- should have been an explicit
deliverable I drafted alongside the v4.1 spec. Instead I wrote the
spec in early May, started Part~2 work on May~20, and had to bolt the
contract on under deadline pressure. The Rev~3 form of
\texttt{part2\_plan.md} works, but it took three drafting rounds to
converge there, two of which were sharpening it after the developer
agent had already started executing.

I would also run Experiment E1 in Part~1, not in Part~2 as a catch-up
entry. The catch-up framing is defensible (entry~17 of my D4 documents
the omission honestly), but the value of E1 is the comparison between
a cold AI grammar and a locked spec, and that comparison would have
been more useful before I locked v4.1 than after -- it might have
surfaced the \texttt{quantity\_of}-as-parens-free issue or the unit-set
drift before the parser was already written.

\paragraph{What did not work and had to be scoped out of v1.}
Recipe composition (decision~\#26) was removed cleanly when I realized
the alternative was building closures and a nested-environment
evaluator that would have doubled the interpreter complexity.
User-defined record types are similarly absent -- the type-name set is
closed in v1 (decision~\#32), which trades extensibility for a sealed
type system that the checker can prove exhaustive. Per-verb action
signatures got collapsed to a three-class table (plan~\S2.3) after I
realized twelve bespoke signatures were $\sim$12 checker branches and
a memorization tax for the exam; the cost is that
\texttt{pour(flour:~Mass)} type-checks even though ``pour''
linguistically implies liquid, and I would address this as a v2
backwards-compatible refinement rather than re-inflate the table now.

The third scoped-out item is more uncomfortable: the type checker does
not flag negative quantities in v1 (spec~\S3, line on unary minus --
interpreter accepts them, type checker does not). This is a real
soundness gap that would matter as soon as a recipe scales by a
negative factor, but I judged that adding a sign discipline to the
dimensional type system was too much work to land safely in the
final hours.

\paragraph{Hardest design decision.} The action-verb three-class
collapse (plan~\S2.3, locking decision~\#29's homogeneous-list
carve-out). It surfaced the underlying design tension between
cooking-faithful semantics (``pour'' should require a liquid,
``blend'' should require same-state ingredients) and type-system
parsimony (one rule per class beats twelve rules per verb). I tested
the decision twice through AI experiments -- once in E3 where Opus~4.7
reproduced the twelve-row mistake unprompted, once in the PM review of
the interpreter commits where the table was probed against the
\S10 error list -- and both confirmed that the three-class form is
the defensible call. But it is still uncomfortable: a Recipix user
who reads ``pour'' in a recipe and finds it accepts a Mass quantity
is going to feel the language doesn't take itself seriously. I made
the call because I have to defend it on the exam in a week and
``twelve places to be wrong'' is harder to defend than ``three places
to be wrong''; v2 can re-impose per-verb dimensions as an optional
refinement without breaking v1 programs.

\vspace{1em}
\noindent\rule{\textwidth}{0.4pt}

\section{Self-Assessment (\textonehalf{} page, six rubric dimensions per handout \S5)}

\paragraph{Correctness -- A.} All three sample programs are wired
end-to-end: samples~1 and~2 produce rendered cookbook-style output,
sample~3 fails at type-check with a clean dimension-mismatch error
message. The 152-test suite passes fully -- zero failures. Sample~1
demonstrates \texttt{scale} correctly multiplying Mass / Volume / Count
ingredients, rebinding the \texttt{servings} parameter, and re-executing
the step body so the \texttt{repeat servings times} loop reflects the
scaled value (six pour/flip cycles for a 4~$\to$~6 servings scale).
Sample~2 demonstrates \texttt{substitute} consuming the
replacement-ingredient binding so the substituted slot reads cleanly;
the non-substituted variant shows the pre-declared replacement
alongside the original, which matches the no-\texttt{this} design
(decision~\#25).

\paragraph{Design Justification -- A$-$.} Every Part-2 contract item
is tied to a written rationale: the three-class action-verb collapse
(plan~\S2.3 + decision~\#29), the in-place AST annotation
(plan~\S2.1), banker's rounding for scaled servings (plan~\S2.6 +
D1~\S4.4), the implicit Ingredient~$\to$~Quantity projection
(decision~\#31 + D1~\S4.8), and the type-equivalence split
(decision~\#10). The minor weakness is that \texttt{scale} skipping
Temperature and Duration is justified in plan~\S4 task~10 with the
cooking-physics argument (intensive vs additive quantities) but isn't
tied to a specific Sebesta hook in the rationale; the professor-agent
review (D4 entry~22) flagged this as a polish-if-time item.

\paragraph{Use of Sebesta Framework -- B.} The major framework moves
are correctly cited: \S6.14 for the structural-vs-name equivalence
split, \S9.5 for the parameter-passing-as-by-value-because-no-mutation
argument, \S7.2 for the precedence and associativity table, \S7.6 for
the left-to-right operand evaluation order. Two places I could tighten:
the \S3 quantity-arithmetic table doesn't cite Sebesta~\S6.2 on
derived/dimensional types for the $Q/Q \to$~unitless rule, and the
spec's by-value paragraph (\S6, decision~\#16) doesn't explicitly name
all four parameter modes from \S9.5 (by-value, by-reference, by-result,
by-value-result) inline. Both are one-sentence fixes I would land if
I had another hour.

\paragraph{Originality -- A$-$.} Recipix's design moves are
domain-specific and defensible against alternatives: the dimensional
type system (decision~\#3), Pinch as its own primitive rather than
\texttt{Quantity<Pinch>} (decision~\#4), \texttt{quantity\_of} as a
unary operator rather than a function because its return type depends
on the operand's quantity-field dimension (decision~\#34),
bare-IDENT substitute slots that make decision~\#28's ``ingredient
identity is the binding'' visible at grammar level (decision~\#35).
E1, E2, and E3 all confirmed that these decisions are not the
textbook defaults -- the cold AI grammars and recommendations missed
each of them.

\paragraph{Code Quality -- A$-$.} The Part-2 modules are cleanly
separated by responsibility: \texttt{runtime\_values.py} is pure data
with one \texttt{normalize\_to\_base} function,
\texttt{environment.py} is type-agnostic shared infrastructure used by
both checker and interpreter, \texttt{interpreter.py} and
\texttt{rendering.py} are Berk-owned and don't touch the type checker,
\texttt{typechecker.py} is İsmail-owned and doesn't touch the
interpreter. No Part-1 contract was modified. The minor code-quality
finding from the professor review is one docstring drift in
\texttt{runtime\_values.py:137} (claims lazy step evaluation;
implementation is eager) which is a 5-line fix in the final pass.

\paragraph{Exam Performance (estimate) -- A$-$.} I can defend the
interpreter mechanics line-by-line: how \texttt{scale} re-executes
step bodies under a rebound \texttt{servings} parameter
(D1~\S4.4) and why that is consistent with referential transparency
(decision~\#25) -- fresh recipe value, no mutation; the banker's
rounding for scaled servings (\texttt{int(round(s * by))}) and its
trade-off against truncation and ceiling; the implicit
Ingredient~$\to$~Quantity projection at every arithmetic call site
(decision~\#31); the loud-raise behavior on \texttt{AmbiguousCall} and
why that matters for integration testing. I can defend the
precedence ladder, the non-associative comparison rule, the
left-to-right operand evaluation order in absence of side effects,
and the assignment-as-statement choice. The one area I am less
rehearsed is articulating spec~\S10 error~\#19 (single-return
enforcement) as a spec-vs-plan reconciliation question --
addressable with a few minutes of review before the 28~May exam.

\end{document}
